\documentclass[12pt]{beamer}

% Tema de la presentación
\usetheme{Madrid} % Puedes cambiarlo: Berlin, CambridgeUS, Warsaw, etc.

% Paquetes recomendados
\usepackage[utf8]{inputenc} % Codificación UTF-8
\usepackage[T1]{fontenc}    % Codificación de fuente
\usepackage[spanish]{babel} % Idioma
\usepackage{graphicx}       % Imágenes
\usepackage{amsmath, amssymb} % Símbolos matemáticos
\usepackage{xcolor}         % Colores personalizados
\usepackage{tikz}
\usetikzlibrary{arrows.meta, positioning}

\newcommand{\lr}{\longrightarrow}


% Datos de la presentación
\title[Mi Presentación]{Título Completo de la Presentación}
\subtitle{Subtítulo opcional}
\author{Tu Nombre}
\institute{Tu Institución}
\date{\today} % o una fecha específica

\begin{document}

% Portada
\begin{frame}
  \titlepage
\end{frame}

% Índice
\begin{frame}{Contenido}
  \tableofcontents
\end{frame}

% Sección 1
\section{Introducción}
\begin{frame}{Introducción}
  El artículo \textbf{"Estimates of earthquake recurrences in the Chaiyi-Tainan area,  Taiwan"} publicado en 2002, aborda un tema de gran revelancia 
  para la sismología y la gestión de riesgos naturales: la estimación de la recurrencias de terremotos en una región altamente sísmica del suroeste de Taiwan.
\end{frame}

\begin{frame}
    \begin{itemize}
        \item Ley de Gutenburg-Richter.
        \item Cadenas de Markov.
        \item Tiene como objetivo estudiar los periodos de  retorno para los terremotos de distinta escala.
        \item Se calcula a partir de la amplitud de las ondas sísmicas registradas en un sismógrafo.
    \end{itemize}
\end{frame}

\begin{frame}{Escala de Richter}
    \begin{itemize}
        \item Es una escala logarítmica utilizada para medir la magnitud de los terremotos.
        \item Fue desarrollada en 1935 por Charles F. Richter
        \item Se basa en la amplitud de las ondas sísmicas.
        \item Cada incremento de 1 en la escala representa un aumento de 10 veces en la amplitud y aproximadamente 31.6 veces más energía liberada.
    \end{itemize}
\end{frame}
% Sección 2
\section{Conceptos Clave}
\begin{frame}
    \begin{block}{Proceso estocástico}
        Es una colección de variables aleatorias que evolucionan en el tiempo o espacio, y que describen un sistema que cambia 
        de manera probabilística. Es un modelo matemático que represneta fenómenos aleatorios que se desarrollan a lo largo de un parámetro,
        esta representado por:
        \[
                \{X_t\}_{t\in T}
        \]
        Donde $X_t$ es una variable aleatoria que representa el estado del sistema en el instante $t$ y $T$ es el conjunto de instantes de tiempo o espacio.
    \end{block}
\end{frame}

\begin{frame}
    \begin{block}{Cadenas de Markov}
        Una cadena de Markov es un tipo especial de proceso estocástico que cumple la propiedad de Markov,
         que establece que el futuro del proceso depende únicamente del estado presente y no de los estados pasados. 
        Formalmente, una cadena de Markov se define por:
        \[
            P(X_{n+1} = x | X_n = x_n, X_{n-1} = x_{n-1}, \ldots, X_0 = x_0) = P(X_{n+1} = x | X_n = x_n)
        \]
        donde $X_n$ es el estado del sistema en el tiempo $n$.
    \end{block}
    Una cadena de Markov posee las siguientes caracteristicas pincipales:
    \begin{itemize}
        \item Propiedad de Markov
        \item Espacio de estados:
        \begin{itemize}
            \item Estados recurentes: aquellos que pueden ser visitados infinitas veces.
            \item Estados transitorios: aquellos que, una vez abandonados, no se vuelven a visitar.
            \item Estados absorbentes: aquellos que, una vez alcanzados, no se abandonan.
        \end{itemize}
        \item Probabilidades de tarnición.
    \end{itemize}
\end{frame}

\begin{frame}{Diagrama de estados}
\begin{figure}[H]
\centering
\begin{tikzpicture}[
    ->, >=Stealth, shorten >=1pt, auto, node distance=3cm,
    state/.style={circle, draw, minimum size=1cm}
]

% Nodos
\node[state] (A) {A};
\node[state] (B) [right=of A] {B};
\node[state] (C) [right=of B] {C};

% Bucles
\path (A) edge[loop above] node{0.2} (A);
\path (C) edge[loop above] node{0.5} (C);

% Flechas entre estados
\path (A) edge[bend left=15] node{0.6} (B)
      (B) edge[bend left=15] node{0.9} (A)
      (B) edge[bend left=15] node{0.5} (C)
      (C) edge[bend left=15] node{0.1} (B)
      (A) edge[bend left=40] node{0.2} (C);

\end{tikzpicture}
\caption{Diagrama de estados con probabilidades de transición.}
\end{figure}
\end{frame}

\begin{frame}{Relación de Chapman-Kolmogorov}
    De acuerdo a la relación de chapman-Kolmogorov suguiere que la matriz de transición 
    $P_{n\times n}$, tiene  la caracteristica de 
    \[
            P^s=P^{s-1}\cdot P
    \]
    Para un estado recurrente, $P_{jj}^{(s)}$ denota la probabilidad de comenzar en el estado $j$
    y regresar a él después de $s$ pasos. Según el teorema del límite para cadenas de Markov, si $j$ es un estado 
    aperiódico, es decir, un estado que se repite en un ciclo no constante,  entonces
    \[
            \lim\limits_{s\lr \infty} p_{jj}^{(s)}=\frac{1}{u_j}
    \]
    Donde $u_j$ es su periodo medio de retorno, que es el número esperado de pasos necesarios para regresar al estado 
    $j$ después de haberlo visitado por primera vez.
\end{frame}

\begin{frame}
    \begin{block}{Definición de matriz ergódica}
        Una matriz de transición $P$ es ergódica si cumple las siguientes
         condiciones:
        \begin{itemize}
            \item Es irreducible: es posible llegar a cualquier estado
             desde cualquier otro estado.
            \item Es aperiódica: no hay ciclos fijos en la cadena de Markov.
            \item Existe un vector propio positivo asociado al valor propio 1,
             que representa la distribución estacionaria.
        \end{itemize}
        En este caso, la cadena de Markov converge a una distribución estacionaria única, independientemente del estado inicial.
        
    \end{block}
\end{frame}

% Sección 3
\section{Conclusiones}
\begin{frame}{Conclusiones}
  \begin{itemize}
    \item Punto importante 1
    \item Punto importante 2
    \item Punto importante 3
  \end{itemize}
\end{frame}

\end{document}
